\documentclass[11pt]{article}
\usepackage[a4paper,margin=22mm]{geometry}
\usepackage{amsmath,amssymb,booktabs,longtable,array,hyperref}
\title{A057 v0.1.1: Master-Mass Reconstruction Blind Falsifier}
\author{Omar Iskandarani}
\date{2026}
\begin{document}\maketitle

% SST-REPORT-SECTION:IDENTIFICATION
\section{Identification}
Catalog ID A057; release v0.1.1; SST Falsifier Framework v1.0.6, profile \texttt{multilibrary\_gpu}. The framework protocol is fail-closed, whereas scientific gates use diagnostic continuation.

% SST-REPORT-SECTION:QUESTION
\section{Question}
Does a topology-stable finite-core vortex energy and its variational spectrum reconstruct a compact mass kernel before protected target comparison?

% SST-REPORT-SECTION:HYPOTHESES
\section{Hypotheses}
The null states that independently sourced carriers do not exhibit a reproducible topology-level energy/dynamical kernel beyond numerical and source variation. The alternative requires reproducible topology structure and protected-holdout generalization before any target reveal.

% SST-REPORT-SECTION:ASSUMPTIONS
\section{Assumptions}
The tested sector is incompressible, inviscid, closed-filament dynamics without reconnection or external forcing. E010/PKLSA is geometry/provenance authority, not particle authority. The physical reference object is the three-dimensional kinetic-energy integral. The operational primary kernel is a regularized slender-filament Hamiltonian; an independent Bishop-frame core-tube quadrature is retained as a second reduced finite-core diagnostic.

% SST-REPORT-SECTION:SYMBOLS
\section{Symbols}
\(\mathbf X_i(s)\) denotes a closed filament component, \(a\) the regularization scale, \(\Gamma\) circulation, \(H_a\) the line Hamiltonian, \(q\) reduced Bishop-frame coordinates, \(K_H\) the energy Hessian and \(J\) the first-order Jacobian.

% SST-REPORT-SECTION:EQUATIONS
\section{Equations}
% AUTO-GENERATED. Do not edit; generated from frozen science_contract.json.
\subsection*{Machine-readable science contract}
\begin{description}
\item[Research question] Does a source-independent, topology-stable finite-core vortex Hamiltonian and its local variational spectrum provide a lower-complexity reconstruction of the historical SST mass kernel, with frozen blind mass-scale predictions that generalize across independent PKLSA carriers?
\item[Objective] Compute geometry, regularized filament energy, interaction coherence, reduced Lagrangian/Hamiltonian closure, Hessian/Jacobian spectra, and target-free surrogate competition on source-native PKLSA carriers. Freeze all topology-level kernels and anonymous composite predictions before any protected target comparison. Scientific gate failures are diagnostic and do not suppress execution of later scientific gates; only protocol/integrity failures are hard blockers.
\item[$H_0$] After numerical convergence and source-independence controls, no registered energy or spectral functional is stable within topology relative to between-topology variation, and added dynamical descriptors do not improve protected holdout prediction of the finite-core energy kernel beyond simple geometry controls.
\item[$H_1$] At least one preregistered finite-core energy/dynamical functional is reproducible across independent carriers of the same topology, separates topology classes beyond numerical/source variation, and yields frozen target-free predictions whose post-reveal mass ratios outperform the registered historical/factorized baselines without target-specific fitting.
\end{description}
\subsubsection*{Registered equations}
\paragraph{E1: Canonical local mass-equivalent density scale}
\begin{equation}
\rho_{\!m}=\frac{\rho_{\!f}\,\|\mathbf{v}_{\!\boldsymbol{\circlearrowleft}}\|^2}{2c^2}
\end{equation}
\noindent\textit{Dimensional check:} (kg m\textasciicircum{}-3)(m\textasciicircum{}2 s\textasciicircum{}-2)/(m\textasciicircum{}2 s\textasciicircum{}-2) = kg m\textasciicircum{}-3.\\
\paragraph{E2: Regularized finite-core filament Hamiltonian}
\begin{equation}
H_a=\frac{\rho\Gamma^2}{8\pi}\sum_{i,j}\oint_{K_i}\!\oint_{K_j}\frac{d\mathbf{X}_i\cdot d\mathbf{X}_j}{\sqrt{\|\mathbf{X}_i-\mathbf{X}_j\|^2+a^2}}
\end{equation}
\noindent\textit{Dimensional check:} rho Gamma\textasciicircum{}2 times a length has units (kg m\textasciicircum{}-3)(m\textasciicircum{}4 s\textasciicircum{}-2)(m)=J.\\
\paragraph{E3: Dimensionless line-energy kernel after length-scale normalization}
\begin{equation}
\widehat{H}_K=\sum_{i,j}\oint\!\oint\frac{d\widehat{\mathbf{X}}_i\cdot d\widehat{\mathbf{X}}_j}{\sqrt{\|\widehat{\mathbf{X}}_i-\widehat{\mathbf{X}}_j\|^2+\widehat a^2}}
\end{equation}
\noindent\textit{Dimensional check:} All hatted lengths are dimensionless, so the complete kernel is dimensionless.\\
\paragraph{E4: Self/mutual coherence matrix}
\begin{equation}
C_{ij}=\frac{E_{ij}}{\sqrt{|E_{ii}E_{jj}|}}
\end{equation}
\noindent\textit{Dimensional check:} Energy divided by energy is dimensionless.\\
\paragraph{E5: Reduced conservative Lagrangian}
\begin{equation}
L(q,\dot q)=\tfrac12\dot q^{\mathsf T}M\dot q-U(q)
\end{equation}
\noindent\textit{Dimensional check:} The blind reduced lane is nondimensionalized; after restoration both terms carry energy units.\\
\paragraph{E6: Reduced Hamiltonian from the Legendre transform}
\begin{equation}
H(q,p)=\tfrac12p^{\mathsf T}M^{-1}p+U(q),\qquad p=M\dot q
\end{equation}
\noindent\textit{Dimensional check:} The blind reduced lane is dimensionless; after restoration H has energy units.\\
\paragraph{E7: Energy Hessian at a registered carrier geometry}
\begin{equation}
K_H=\left.\nabla_q^2U(q)\right|_{q=0}
\end{equation}
\noindent\textit{Dimensional check:} Dimensionless in the normalized lane.\\
\paragraph{E8: First-order dynamical Jacobian}
\begin{equation}
J=\begin{pmatrix}0&I\\-M^{-1}K_H&-M^{-1}G\end{pmatrix}
\end{equation}
\noindent\textit{Dimensional check:} In normalized time J is dimensionless per unit normalized time; a chosen physical time scale restores s\textasciicircum{}-1.\\
\paragraph{E9: Spectral ratio with no protected target embedded}
\begin{equation}
q_{21}(K)=\omega_2(K)/\omega_1(K)
\end{equation}
\noindent\textit{Dimensional check:} Ratio of equal-frequency units is dimensionless.\\
\paragraph{E10: SI mass prediction from the regularized Hamiltonian}
\begin{equation}
M_K^{(\rho)}=H_K^{(\rho)}/c^2
\end{equation}
\noindent\textit{Dimensional check:} J/(m\textasciicircum{}2 s\textasciicircum{}-2)=kg.\\
\paragraph{E11: Three-dimensional kinetic-energy mass functional used as the physical reference object}
\begin{equation}
M_K^{(3D)}=\frac{1}{2c^2}\int_{\mathbb R^3}\rho_K(\mathbf x)\|\mathbf v_K(\mathbf x)\|^2\,d^3x
\end{equation}
\noindent\textit{Dimensional check:} (kg m\textasciicircum{}-3)(m\textasciicircum{}2 s\textasciicircum{}-2)(m\textasciicircum{}3)/(m\textasciicircum{}2 s\textasciicircum{}-2)=kg.\\
\paragraph{E12: Core-tube quadrature kernel used as an independent reduced check on the line Hamiltonian}
\begin{equation}
\widehat E_{\mathrm{tube}}=\frac12\int_{\widehat\Omega_{\mathrm{tube}}}\|\widehat{\mathbf v}(\widehat{\mathbf x})\|^2\,d^3\widehat x,\qquad E_{\mathrm{tube}}=\frac{\rho\Gamma^2r_c}{32\pi^2}\widehat E_{\mathrm{tube}}
\end{equation}
\noindent\textit{Dimensional check:} The hatted integral is dimensionless; rho Gamma\textasciicircum{}2 r\_c has units J.\\
\subsubsection*{Registered code-to-mathematics steps}
\begin{enumerate}
\item \textbf{S1} [G1; equations E3]: Verify E010 release metadata, locate qualification ledgers, reload original carriers, and recheck raw/geometry hashes when available.
\item \textbf{S2} [G2; equations E1, E2, E3]: Run analytic/synthetic controls, SI scale checks, orientation/rigid-motion/resampling controls.
\item \textbf{S3} [G3; equations E3]: Compute length, reach-normalized ropelength proxy, curvature, torsion, writhe/linking, Bishop holonomy, and convergence statistics.
\item \textbf{S4} [G4; equations E3]: Evaluate held-out independent carrier groups without pooling distinct carriers as numerical refinement levels.
\item \textbf{S5} [G5; equations E2, E3]: Compare C++/pybind11 FP64 regularized energy matrix with Python/NumPy FP64 reference.
\item \textbf{S6} [G6; equations E2]: Record framework SYCL/DD32 screening evidence separately from certification authority.
\item \textbf{S7} [G7; equations E3]: Quantify cross-independence-family replication and pseudo-replication exclusions.
\item \textbf{S8} [G8; equations E2, E4, E11, E12]: Compute self/mutual energy matrices, coherence spectra, and an independent Bishop-frame core-tube kinetic-energy quadrature for a bounded topology subset.
\item \textbf{S9} [G10; equations E5]: Evaluate reduced Euler-Lagrange residuals around source-native geometries even when earlier source/geometry gates fail, carrying caveats forward.
\item \textbf{S10} [G11; equations E5, E6]: Evaluate Legendre-transform closure and reduced Hamiltonian conservation.
\item \textbf{S11} [G12; equations E7]: Construct finite-difference Hessians under Bishop-frame perturbations and test symmetry/refinement.
\item \textbf{S12} [G13; equations E8, E9]: Construct Jacobians, extract spectra, and compare oscillatory frequencies with direct reduced integration.
\item \textbf{S13} [G14; equations E9]: Measure topology stability of spectral ratios without accessing the protected historical target.
\item \textbf{S14} [G15; equations E3]: Cross-validate simple geometry surrogates against the finite-core energy kernel with topology-grouped holdouts.
\item \textbf{S15} [G16; equations E3, E9]: Compare preregistered energy-only and geometry-plus-spectrum surrogate families under a complexity penalty.
\item \textbf{S16} [G17; equations E3, E10]: Freeze a complete target-free topology prediction table and anonymous composite predictions before reveal.
\item \textbf{S17} [G18; equations E3]: Evaluate protected topology holdout generalization internal to the geometry/energy task; no observed mass targets are used.
\item \textbf{S18} [G19; equations E2, E10]: Restore SI scales using canonical SST constants in two density lanes and write target-free absolute predictions.
\item \textbf{S19} [G20; equations E2, E7, E8]: Stress core regularization, discretization, carrier selection, and perturbation amplitude; quantify sensitivity.
\item \textbf{S20} [G21; equations E3, E9]: Jackknife across provider/lineage groups and report whether findings replicate in independent evidence families.
\item \textbf{S21} [G22; equations E3, E5, E6, E8, E10]: Assign a blind multi-level classification from all recorded evidence while preserving diagnostic results from failed upstream scientific gates.
\end{enumerate}


% SST-REPORT-SECTION:ALGORITHM
\section{Algorithm}
Carriers are reloaded from source-native bytes, independently resampled, reach-normalized, and evaluated on preregistered resolution/core ladders. Geometry, self/mutual line energy, an independent core-tube kinetic-energy quadrature, coherence, reduced-action residuals, Hessians and Jacobians are computed without protected targets. Topology-grouped model competition predicts the energy kernel rather than observed masses.

% SST-REPORT-SECTION:SOURCES
\section{Sources}
The primary runtime source is E010 PKLSA v0.3.x. Carrier-level hash verification and independence labels are rechecked at runtime. Synthetic controls are implementation evidence only.

% SST-REPORT-SECTION:GATES
\section{Gates and diagnostic continuation}
Only G0 is a hard framework prerequisite. Every scientific gate depends directly on G0 rather than on the preceding scientific gate. Consequently a scientific FAIL does not prevent later computation. A separate soft-dependency ledger records interpretive caveats. This distinction is intentional: evidence is retained while unsupported promotion is prevented.

% SST-REPORT-SECTION:NUMERICS
\section{Numerics}
BASIC uses a small selected set; FULL scans admitted topology classes through the registered crossing limit; CERTIFY increases carrier and resolution ladders. Regularized double line integrals are evaluated in FP64. A separate tube-support velocity quadrature probes the three-dimensional kinetic-energy route on a bounded subset. Hessians use centered finite differences under small Bishop-frame perturbations.

% SST-REPORT-SECTION:BACKENDS
\section{Backends}
Python/NumPy FP64 is REFERENCE. C++/pybind11 FP64 is CERTIFICATION and must satisfy relative-L2 \(\le10^{-10}\). SYCL FP32 and DD32/FP32x2 remain SCREENING\_ONLY under framework v1.0.6 and cannot certify the scientific result.

% SST-REPORT-SECTION:BLINDNESS
\section{Blindness}
Protected observed targets, protected human identities and protected historical numerical comparison values reside only under \texttt{private/}. The blind run freezes a complete topology prediction table first. Reveal requires the v1.0.6 output-integrity and nonced-commitment checks.

% SST-REPORT-SECTION:STATISTICS
\section{Statistics}
Primary diagnostics are within-topology coefficient of variation, between/within spread ratio, independent-group replication fraction, topology-grouped confirmation/holdout log-RMSE, carrier jackknife sensitivity, core-law sensitivity and target-free spectral-ratio stability.

% SST-REPORT-SECTION:RESULTS
\section{Results}
\IfFileExists{AUTO_RESULTS.tex}{\input{AUTO_RESULTS.tex}}{No run results are embedded in the source release.}

% SST-REPORT-SECTION:INTERPRETATION
\section{Interpretation}
A later diagnostic may remain numerically informative even when an earlier scientific gate fails. Such a result is explicitly downgraded by \texttt{DEPENDENCY\_CAVEATS.json}; for example a Jacobian around a nonstationary carrier is a local numerical diagnostic rather than a certified stability spectrum. A mass claim requires the separately registered blind prediction and post-reveal comparison.

% SST-REPORT-SECTION:REPRODUCIBILITY
\section{Reproducibility}
Run \texttt{run\_00\_install.cmd}, verify/freeze the protocol, then execute BASIC, FULL or CERTIFY. The canonical Workbench root defaults to \texttt{C:\textbackslash workspace\textbackslash projects\textbackslash SST-Workbench} and may be overridden by \texttt{SST\_WORKBENCH\_ROOT}.

\begin{thebibliography}{99}
\bibitem{Moffatt1969} H.~K. Moffatt, 1969, ``The degree of knottedness of tangled vortex lines,'' \emph{Journal of Fluid Mechanics} 35, 117--129. DOI: \href{https://doi.org/10.1017/S0022112069000991}{10.1017/S0022112069000991}.
\bibitem{Saffman1992} P.~G. Saffman, 1992, \emph{Vortex Dynamics}, Cambridge University Press. DOI: \href{https://doi.org/10.1017/CBO9780511624063}{10.1017/CBO9780511624063}.
\bibitem{Bishop1975} R.~L. Bishop, 1975, ``There is More than One Way to Frame a Curve,'' \emph{American Mathematical Monthly} 82, 246--251. DOI: \href{https://doi.org/10.1080/00029890.1975.11993807}{10.1080/00029890.1975.11993807}.
\bibitem{Adams1991} C.~C. Adams, M.~Hildebrand, and J.~Weeks, 1991, ``Hyperbolic invariants of knots and links,'' \emph{Transactions of the American Mathematical Society} 326, 1--56. DOI: \href{https://doi.org/10.1090/S0002-9947-1991-0994161-2}{10.1090/S0002-9947-1991-0994161-2}.
\bibitem{Helmholtz1858} H.~von Helmholtz, 1858, ``\"Uber Integrale der hydrodynamischen Gleichungen, welche den Wirbelbewegungen entsprechen,'' \emph{Journal f\"ur die reine und angewandte Mathematik} 55, 25--55. Permalink: \href{https://eudml.org/doc/147531}{eudml.org/doc/147531}.
\end{thebibliography}
\end{document}
